% Appendix A (track hw). The former Background section of the main text (the five files
% "bg - *.tex" of the 438f6f16 mirror, coauthor text), moved here, deduplicated against
% Secs. 2.1 to 2.6 and trimmed. The gate actions are Eq. (eq:conventions) of Sec. 2.1; the
% paper's own protocol choices (mirror inverse, estimator, zero offset, unit of depth) are
% Secs. 2.3 and 2.6 and are not restated here. Changes of content are listed in
% jpw-review/hw/content-changes.md, part F. The "bg - *.tex" files stay in the tree, unused.
\section{Background}
\label{app:background}\label{sec:bg}

This appendix collects the standard material that the Methods build on: the qudit Pauli and Clifford groups with the conjugation rules of the gates of Eq.~\eqref{eq:conventions}, randomized benchmarking and its logical form, and the qudit surface code with its folded variant.

\subsection{Qudit Pauli and Clifford Groups}
\label{app:background-groups}

% from "bg - Qudit Pauli Group.tex", lines 2-36
The Pauli group of a qudit of prime dimension $q$ is generated by the operators $X$ and $Z$ of Eq.~\eqref{eq:conventions} together with the phase $\omega I$, where $I$ is the identity~\cite{Gottesman1999,Hostens2005,Wang2020,SarkarYoder2024},
\begin{equation}
  \mathcal{P}_q = \bigl\{\, \omega^c X^a Z^b : a, b, c \in \mathbb{Z}_q \,\bigr\},
\end{equation}
with $\mathbb{Z}_q$ the integers modulo $q$. The generators obey $X^q = Z^q = I$ and $Z X = \omega\, X Z$, so that $X^a Z^b = \omega^{-ab}\, Z^b X^a$, and every element of $\mathcal{P}_q$ has a unique triple $(a, b, c)$. The Pauli operators form a discrete basis for arbitrary errors on the qudit. For $n$ qudits the Pauli group is
\begin{equation}
  \mathcal{P}_q^{\otimes n}
  = \bigl\{\, \omega^c
      X_1^{a_1} \cdots X_n^{a_n}\,
      Z_1^{b_1} \cdots Z_n^{b_n}
      : a_1, \dots, a_n, b_1, \dots, b_n, c \in \mathbb{Z}_q
     \,\bigr\}.
\end{equation}

% from "bg - Qudit Clifford Group.tex", lines 3-57
The $n$-qudit Clifford group is the normalizer of $\mathcal{P}_q^{\otimes n}$ in the unitary group $U(q^n)$, that is, the set of unitaries $U$ with $U \mathcal{P}_q^{\otimes n} U^\dagger = \mathcal{P}_q^{\otimes n}$~\cite{Gottesman1999}. We write $\mathcal{C}_q$ for the single-qudit Clifford group, taken up to global phases. Circuits that consist of Clifford gates and Pauli measurements can be simulated efficiently on a classical computer~\cite{Gottesman1998}, which is the basis of the stabilizer simulations of \autoref{sub:sim_framework}. For prime $q$ the Clifford group is generated by $H$, $S$ and SUM~\cite{Hostens2005}. These gates conjugate the Pauli operators as
\begin{equation}
  \begin{aligned}
    H X H^\dagger &= Z, &\qquad H Z H^\dagger &= X^{-1}, \\
    S X S^\dagger &= X Z, &\qquad S Z S^\dagger &= Z, \\
    \mathrm{SUM}\,(X \otimes I)\,\mathrm{SUM}^\dagger &= X \otimes X, &\qquad \mathrm{SUM}\,(I \otimes X)\,\mathrm{SUM}^\dagger &= I \otimes X, \\
    \mathrm{SUM}\,(Z \otimes I)\,\mathrm{SUM}^\dagger &= Z \otimes I, &\qquad \mathrm{SUM}\,(I \otimes Z)\,\mathrm{SUM}^\dagger &= Z^{-1} \otimes Z,
  \end{aligned}
  \label{eq:bg-conjugation}
\end{equation}
where the first tensor factor is the control of SUM. The gate $H$ is the discrete Fourier transform over $\mathbb{Z}_q$ and exchanges the eigenbases of $Z$ and $X$. The gate $S$ is $\mathrm{diag}(1, 1, \omega)$ for $q = 3$. The factor $Z^{-1}$ that SUM leaves on its control has no counterpart for qubits, where $Z^{-1} = Z$. The gate CZ follows from the generators as $\mathrm{CZ} = (I \otimes H)\,\mathrm{SUM}\,(I \otimes H^\dagger)$. It commutes with $Z \otimes I$ and $I \otimes Z$ and maps $X \otimes I$ to $X \otimes Z$ and $I \otimes X$ to $Z \otimes X$.

\subsection{Randomized Benchmarking}
\label{app:background-rb}

% from "bg - RB.tex", lines 2-19, 59-65 and 88
Randomized benchmarking (RB) estimates the average error rate of a set of gates from the decay of a signal under random gate sequences~\cite{Emerson2005,Knill2008,Magesan2011PRL}. It was developed for qubits and carries over to qudits, because the Clifford group of a qudit of prime dimension is also a unitary 2-design~\cite{Appleby2005,Dankert2009,Gross2007}. The standard protocol for one qutrit~\cite{Appleby2005,Morvan2021} prepares $\ket{0}$, applies a sequence of $m$ Cliffords $C_1, \dots, C_m$ drawn uniformly from $\mathcal{C}_3$, applies the single recovery gate $C_{\mathrm{inv}} = (C_m \cdots C_1)^{-1}$, which returns the qutrit to $\ket{0}$ in the absence of errors, and measures in the computational basis. Averaged over many sequences, the probability of the outcome zero depends on the sequence length as
\begin{equation}
F_{\mathrm{seq}}(m) = A f^{m} + B
\label{eq:fseq}
\end{equation}
if the noise is Markovian and does not depend on the gate~\cite{Magesan2011PRL,Magesan2012PRA}. The decay parameter $f$ characterizes the gates, and the constants $A$ and $B$ absorb the errors of state preparation and measurement.

% from "bg - RB.tex", lines 21-56
The single exponential of Eq.~\eqref{eq:fseq} follows from the 2-design property. Let $\Lambda$ be a completely positive trace-preserving noise channel that follows each Clifford. Its average over conjugation by the Clifford group, the twirl
\begin{equation}
    \Lambda_{\mathcal{C}_3}(\rho)
    = \frac{1}{|\mathcal{C}_3|}\sum_{C\in\mathcal{C}_3}
    C^\dagger\,\Lambda(C\rho\, C^\dagger)\,C
    = \int_{U(3)} U^\dagger\,\Lambda(U\rho\, U^\dagger)\,U\,
    d\mu(U),
\label{eq:clifford-twirl}
\end{equation}
equals the average over the Haar measure $\mu$ of the unitary group and is therefore a depolarizing channel with the average fidelity of $\Lambda$~\cite{Emerson2005,Dankert2009,Gross2007,Magesan2011PRL},
\begin{equation}
    \Lambda_{\mathcal{C}_3}(\rho) = (1-p_{\mathrm{mix}})\,\rho+\frac{p_{\mathrm{mix}}}{3}\,I ,
    \label{eq:1-depol}
\end{equation}
which replaces the state by the completely mixed state with probability $p_{\mathrm{mix}}$~\cite{NielsenChuang2010}. In the language of \autoref{sec:rep-theory}, conjugation by the Clifford group splits the operators on the qutrit into two irreducible subspaces, the span of the identity and the span of the $q^2 - 1 = 8$ traceless Pauli operators. By Schur's lemma the twirled channel acts on each of them as a multiple of the identity,
\begin{equation}
    \Lambda_{\mathcal{C}_3} =
    \begin{bmatrix}
    1 & 0 \\
    0 & f\,\Pi
    \end{bmatrix},
    \qquad
    \Lambda_{\mathcal{C}_3}^{\,m} =
    \begin{bmatrix}
    1 & 0 \\
    0 & f^{m}\,\Pi
    \end{bmatrix},
\label{eq:twirl_irrep}
\end{equation}
where $\Pi$ is the identity on the traceless subspace and $f = 1 - p_{\mathrm{mix}}$. The twirled channel thus has the two eigenvalues $1$ and $f$. The identity block gives the constant $B$ of Eq.~\eqref{eq:fseq} and the traceless block the term $A f^m$.

% from "bg - RB.tex", lines 2-7 and 88-94
The decay parameter gives the average gate fidelity~\cite{Nielsen2002,Magesan2012PRA},
\begin{equation}
\bar{F} = \frac{(q-1)f + 1}{q} ,
\label{eq:fbar}
\end{equation}
of which Eq.~\eqref{eq:error_rate} is the case $q = 3$. When the noise depends weakly on the gate, the decay stays close to a single exponential, whereas drift, non-Markovian noise and leakage can break this form~\cite{Emerson2005,Wallman2018}. Within these limits, RB with Clifford gates and depolarizing noise of a suitably chosen probability reproduces the decay of hardware with more general errors~\cite{Magesan2011PRL,Proctor2019DRB}, and the method is established for qutrits~\cite{Morvan2021}. \autoref{fig:rb-stacked} shows the mirror variant of the RB circuit that serves as our baseline (\autoref{sub:rb_baseline}).

\begin{figure}[tb]
  \centering
    \begin{quantikz}[column sep=0.22cm]
      \lstick{\textcolor{LightBlue}{$\ket{0}$}} &
        \gate{\textcolor{red}{C_1}} &
        \gate{\textcolor{violet}{\mathcal{N}}} &
        \push{\,\cdots\,} &
        \gate{\textcolor{red}{C_m}} &
        \gate{\textcolor{violet}{\mathcal{N}}} &
        \gate{\textcolor{red}{C_m^{-1}}} &
        \gate{\textcolor{violet}{\mathcal{N}}} &
        \push{\,\cdots\,} &
        \gate{\textcolor{red}{C_1^{-1}}} &
        \gate{\textcolor{violet}{\mathcal{N}}} &
        \meter{}
    \end{quantikz}
  \caption{RB circuit of depth $m$ in the mirror variant that serves as the baseline (\autoref{sub:rb_baseline}). The qutrit starts in $\ket{0}$, the Cliffords $C_1, \dots, C_m$ are followed by their inverses in reverse order, and the qutrit is measured in the computational basis. Each box $\mathcal{N}$ stands for the noise that follows the primitive gates of the preceding Clifford (\autoref{sub:noise_models}). The standard protocol replaces the inverse half by the single recovery gate $C_{\mathrm{inv}}$.}
  \label{fig:rb-stacked}
\end{figure}

\subsection{Logical Randomized Benchmarking}
\label{app:background-lrb}

% from "bg - LRB.tex", lines 3-43
Logical randomized benchmarking, proposed by Combes \emph{et al.}~\cite{Combes2017}, applies RB to the logical gates of an error-correcting code, so that the error rate at the logical level can be measured and compared with the physical one. The protocol prepares a logical basis state such as $\ket{\bar{0}}$ or $\ket{\bar{+}}$ and applies $m$ random logical Cliffords $\bar C_1, \dots, \bar C_m$, each realized by physical operations on the qudits of the code. A round of error correction follows each logical Clifford: it measures the stabilizers and applies a recovery that returns the state to the code space. The single logical inverse $\bar C_{\mathrm{inv}} = (\bar C_m \cdots \bar C_1)^{-1}$ and the measurement of a logical operator end the circuit. The probability that this measurement agrees with the initial logical state is fitted over $m$ as in RB, and its decay parameter gives the average logical error rate per Clifford~\cite{Magesan2011PRL,Magesan2012PRA,Wallman2018}.

% from "bg - LRB.tex", line 46
The protocol characterizes an implementation of error correction at the logical level directly. It avoids the extrapolation from physical error rates, which can misestimate the logical performance by orders of magnitude. Because it works in the logical code space, it needs fewer experimental resources than a characterization of the physical level, and it estimates the average probabilities of correctable and uncorrectable errors and the average error rate of the logical gates. It includes effects that physical-level models often neglect, namely imperfect logical gates, crosstalk and correlated errors. Its proof of correctness reduces to that of RB, so that it inherits the robustness and the analysis tools of RB~\cite{Combes2017}.

% from "bg - LRB.tex", lines 43 and 48-51
These guarantees rest on the assumptions of RB, noise on the logical Cliffords that is time-independent, Markovian and gate-independent, and on the additional assumption that syndrome measurement errors are uncorrelated between rounds of error correction. The protocol tolerates small violations of these assumptions, whereas large ones compromise the interpretation of the fitted parameters. It also requires a code with implemented logical Cliffords and syndrome readout, so it does not replace the characterization of the physical hardware~\cite{Combes2017}.

\subsection{Qudit Surface Code and Folding}
\label{sec: qudit surface code}\label{app:background-folded}

% from "bg - Modified Qutrit Surface Code.tex", line 3
The surface code generalizes from qubits to qudits of dimension $q$ when qudit Pauli operators serve as stabilizers. The qudits sit on the edges of a directed planar graph (\autoref{fig:qudit-surface-code}). Each vertex carries an $X$-type stabilizer generator and each face a $Z$-type generator, products of $X$ or of $Z$ and their inverses on the adjacent edges, in which the edge directions decide where the inverses stand. These generators commute and define a code space that encodes one logical qudit.
% ============================================================
% fig - qudit surface code.tex
% Figure: qudit surface code stabilizers and distance-5 patch.
% Shared colors and sizes come from 'tikz preamble.tex'.
% All three panels are drawn at the same scale (\SCscale) with a
% unit lattice cell, so stabilizer close-ups and the full patch
% stay proportional.
% ============================================================

% ---------- local helpers for this figure only ----------
\newcommand{\SCscale}{1.5}
\newcommand{\SCctx}{black!30}          % context lattice edges in close-ups
\newcommand{\SCgrid}{0.9pt}            % lattice line width
\newcommand{\SCctxline}{0.6pt}          % context lattice line width
\newcommand{\SCpatchline}{1.0pt}       % stabilizer patch outline width

% One qudit: white circle + orientation wedge (0 = right, 90 = up).
\newcommand{\SCqudit}[3]{% position, wedge angle, label
  \WedgeArrow{#1}{#2}
  \draw[line width=\qdLine,fill=white] (#1) circle (\qdR);
  \node[font=\numFont] at (#1) {#3};
}

\begin{figure}[H]
  \centering
  \captionsetup[subfigure]{justification=centering}
  % Left column: the two stabilizer close-ups, drawn on a 2x2-cell window.
  \begin{minipage}[t]{0.34\textwidth}
    \vspace{0pt}
    % -------- (a) Vertex / X-type check --------
    \begin{subfigure}[t]{\linewidth}
      \centering
      \begin{tikzpicture}[line cap=round,line join=round,scale=\SCscale]
        % context lattice
        \foreach \y in {0,1,2} \draw[\SCctx,line width=\SCctxline] (0,\y) -- (2,\y);
        \foreach \x in {0,1,2} \draw[\SCctx,line width=\SCctxline] (\x,0) -- (\x,2);
        % stabilizer patch: diamond over the four edge midpoints
        \path[fill=Xvivid,draw=Xcol!55!black,line width=\SCpatchline]
          (1,1.5) -- (1.5,1) -- (1,0.5) -- (0.5,1) -- cycle;
        % the four edges of the central vertex, at full strength
        \draw[line width=\SCgrid] (0,1) -- (2,1);
        \draw[line width=\SCgrid] (1,0) -- (1,2);
        % qudits
        \SCqudit{1,1.5}{90}{n}
        \SCqudit{1.5,1}{0}{e}
        \SCqudit{1,0.5}{90}{s}
        \SCqudit{0.5,1}{0}{w}
      \end{tikzpicture}
      \caption{Vertex stabilizer generator \mbox{$G_X=X^{\dagger}_n \otimes X_w \otimes X^{\dagger}_e \otimes X_s$}.}
      \label{fig:vertex}
    \end{subfigure}

    \vspace{1.2ex}

    % -------- (b) Face / Z-type check --------
    \begin{subfigure}[t]{\linewidth}
      \centering
      \begin{tikzpicture}[line cap=round,line join=round,scale=\SCscale]
        % context lattice: neighboring edges plus the window frame
        \foreach \y in {0.5,1.5} \draw[\SCctx,line width=\SCctxline] (0,\y) -- (2,\y);
        \foreach \x in {0.5,1.5} \draw[\SCctx,line width=\SCctxline] (\x,0) -- (\x,2);
        \draw[\SCctx,line width=\SCctxline] (0,0) rectangle (2,2);
        % stabilizer patch: diamond over the four edge midpoints of the face
        \path[fill=Zvivid,draw=Zcol!60!black,line width=\SCpatchline]
          (1,1.5) -- (1.5,1) -- (1,0.5) -- (0.5,1) -- cycle;
        % the face itself, at full strength
        \draw[line width=\SCgrid] (0.5,0.5) rectangle (1.5,1.5);
        % qudits
        \SCqudit{1,1.5}{0}{n}
        \SCqudit{1.5,1}{90}{e}
        \SCqudit{1,0.5}{0}{s}
        \SCqudit{0.5,1}{90}{w}
      \end{tikzpicture}
      \caption{Face stabilizer generator \mbox{$G_Z = Z_n \otimes Z_w \otimes Z^{\dagger}_e \otimes Z^{\dagger}_s$}.}
      \label{fig:face}
    \end{subfigure}
  \end{minipage}
  \hfill
  % Right column: the distance-5 patch at the same scale.
  \begin{minipage}[t]{0.62\textwidth}
    \vspace{0pt}
    \begin{subfigure}[t]{\linewidth}
      \centering
      \begin{tikzpicture}[line cap=round,line join=round,scale=\SCscale]
        \newcommand{\nCols}{5}
        \newcommand{\nRows}{5}
        \newcommand{\xLabelOffset}{0.52}
        \newcommand{\zLabelDy}{0.42}

        % top X-type boundary triangles
        \foreach \i in {0,...,\numexpr\nCols-2\relax} {
          \path[fill=Xvivid, draw=none]
            (\i+0.5,\nRows) -- (\i+1,\nRows-0.5) -- (\i+1.5,\nRows) -- cycle;
        }
        % interior X-type diamonds
        \foreach \x in {1,...,\numexpr\nCols-1\relax} {
          \foreach \y in {2,...,\numexpr\nRows-1\relax} {
            \path[fill=Xvivid, draw=none]
              (\x,\y+0.5) -- (\x+0.5,\y) -- (\x,\y-0.5) -- (\x-0.5,\y) -- cycle;
          }
        }
        % Z-type cells: left triangles, interior diamonds, right triangles
        \foreach \y in {2,...,\numexpr\nRows\relax} {
          \path[fill=Zvivid, draw=none]
            (0.5,\y) -- (1,\y-0.5) -- (0.5,\y-1) -- cycle;
        }
        \foreach \x in {1,...,\numexpr\nCols-2\relax} {
          \foreach \y in {2,...,\numexpr\nRows\relax} {
            \path[fill=Zvivid, draw=none]
              (\x+0.5,\y) -- (\x+1,\y-0.5) -- (\x+0.5,\y-1) -- (\x,\y-0.5) -- cycle;
          }
        }
        \foreach \y in {2,...,\numexpr\nRows\relax} {
          \path[fill=Zvivid, draw=none]
            (\nCols-0.5,\y) -- (\nCols-1,\y-0.5) -- (\nCols-0.5,\y-1) -- cycle;
        }

        % bottom X-type boundary triangles
        \foreach \i in {0,...,\numexpr\nCols-2\relax} {
          \path[fill=Xvivid, draw=none]
            (\i+0.5,1) -- (\i+1,1.5) -- (\i+1.5,1) -- cycle;
        }

        % lattice
        \foreach \y in {1,...,\nRows}
          \draw[line width=\SCgrid] (0.5,\y)--(\nCols-0.5,\y);
        \foreach \x in {1,...,\numexpr\nCols-1\relax}
          \draw[line width=\SCgrid] (\x,1)--(\x,\nRows);

        % qudits with orientation wedges (right on horizontal edges, up on vertical)
        \newcounter{scqid}\setcounter{scqid}{0}
        \foreach \j in {0,...,8} {
          \pgfmathsetmacro\yy{\nRows - 0.5*\j}
          \ifodd\j\relax
            \foreach \x in {1,...,\numexpr\nCols-1\relax} {
              \stepcounter{scqid}
              \SCqudit{\x,\yy}{90}{\arabic{scqid}}
            }
          \else
            \foreach \x in {0,...,\numexpr\nCols-1\relax} {
              \stepcounter{scqid}
              \SCqudit{\x+0.5,\yy}{0}{\arabic{scqid}}
            }
          \fi
        }

        % boundary labels
        \foreach \j in {0,2,4,6,8} {
          \pgfmathsetmacro\yy{\nRows - 0.5*\j}
          \node[text=Xcol!80!black] at (0.5-\xLabelOffset,\yy) {\labelFont $X$};
        }
        \foreach \x in {0,...,\numexpr\nCols-1\relax} {
          \node[text=Zcol!85!black] at (\x+0.5,\nRows+\zLabelDy) {\labelFont $Z$};
        }
      \end{tikzpicture}
      \caption{Distance-5 patch.}
      \label{fig:patch}
    \end{subfigure}
  \end{minipage}

  % Caption revised by track hw (2026-11-13) for Appendix A. The former caption is in paper-rework/draft/.
  % TODO(authors): the former caption said "(a) X acts on outgoing edges, X^dagger on incoming edges". With the edge directions drawn in panel (a) (horizontal edges point right, vertical edges point up), the formula of the sub-caption, G_X = X_n^dagger X_w X_e^dagger X_s, puts X^dagger on the outgoing edges n and e. The caption below states the rule without the assignment; please confirm which is meant.
  \caption{Qudit surface code~\cite{Moussa2016}. Circles are qudits on the edges of a directed planar graph, and the wedge at a circle points along its edge. Blue marks $X$-type and red $Z$-type stabilizer generators.
  (a)~Vertex generator $G_X$: $X$ or $X^\dagger$ on each edge at the vertex, set by whether the edge points into the vertex or out of it.
  (b)~Face generator $G_Z$: $Z$ on the clockwise edges of the face and $Z^\dagger$ on the counterclockwise ones.
  (c)~Distance-5 patch of 41 qudits. Truncated $X$-type generators close the top and bottom boundaries and truncated $Z$-type generators the left and right ones. The labels $X$ and $Z$ mark the supports of the logical operators: $\bar X$ acts on the left column, between the top and bottom boundaries, and $\bar Z$ on the top row, between the left and right boundaries.}
  \label{fig:qudit-surface-code}
\end{figure}


% from "bg - Modified Qutrit Surface Code.tex", lines 6-12
% TODO(authors): the former text wrote the gates as formulas in a notation that was not defined, Ubar = (U[v_bullet] x U^dagger[v_star])(U'[v_i,bullet] x U'[v_i,star]) with SWAP and CZ on the pairs. Moussa's paper states the gates in words (H and H^dagger on the fold qudits and a mirror gate on the pairs; S and S^dagger on the fold qudits and CZ and CZ^dagger on the pairs), and the paragraph below follows that statement. Please confirm.
Moussa~\cite{Moussa2016} introduced folded surface codes to give surface codes the constant-depth transversal single-qudit Clifford gates of color codes. A surface code by itself lacks the self-dual symmetry that makes the Clifford gates of a color code transversal, because each of its generators is of one type only. Folding the patch supplies a weaker, mirror-dual symmetry: it pairs each $X$-type generator of one layer with a $Z$-type generator of the other. The fold places pairs of qudits on top of each other, and the qudits on the fold line stay single. A bipartition of the qudits into two types, the star bipartition, decides which qudits carry a gate and which its inverse, and the fold only pairs qudits of the same type. The logical $\bar H$ then applies $H$ or $H^\dagger$ to every qudit and a SWAP to every pair, and the logical $\bar S$ applies $S$ or $S^\dagger$ to the qudits on the fold line and CZ or $\mathrm{CZ}^\dagger$ to every pair. Equations~\eqref{eq:logical_H} and~\eqref{eq:logical_HS} are these gates for the five-qutrit patch.

On planar hardware, two qudits that the fold pairs can lie far apart. The folding step of~\cite{Moussa2016} brings them together: SWAP gates between data qudits and additional ancilla qudits move the two layers on top of each other, so that the gates on a pair become local. The folding step acts as the identity on the logical qudit and changes only the layout, and its circuit depth grows linearly with the code distance. The construction builds on the equivalence between color codes and folded surface codes~\cite{Kubica2015} and has been extended to further Calderbank--Shor--Steane codes through fold-transversal gates~\cite{BreuckmannBurton2024}.
